\section{Where the work lies}

The analysis assesses the workflow task by task rather than as a whole.
Studies of how technology changes work divide jobs into tasks and ask which of them can be automated, because a job is rarely automated as a whole \citep{autor2003skill}.
\citet{eloundou2024gpts} applied this to LLMs, rating each task of an occupation by whether access to a model alone, or to a model with additional software, could at least halve the time it takes at the same quality.
Judging tasks one at a time is also a safeguard: in a field experiment with consultants, AI improved work on tasks within its capabilities but made it worse on a task just outside them that looked similar \citep{dellacqua2026navigating}.
Within each step, the four questions of Section 4.5 decide where each part belongs: whether a rule can compute one correct answer, whether the input is unstructured, whether the result can be checked cheaply, and whether an error would be costly and hard to see.

Table \ref{tab:where-work-lies} gives the result for each step, and Figure \ref{fig:task-map} places the main parts on two of those questions: how much judgment a part needs, and how quietly an error in it would pass.

\begin{table}[htbp]
\centering
\small
\caption{What each step of the workflow needs, and where its parts belong.}
\label{tab:where-work-lies}
\newcommand{\hdr}[2]{\textcolor{#1}{\rule{1.8mm}{1.8mm}}\enspace\textbf{#2}}
\setlength{\tabcolsep}{5pt}
\begin{tabularx}{\textwidth}{@{}>{\raggedright\arraybackslash\bfseries}p{2.1cm}*{3}{>{\raggedright\arraybackslash\setlength{\parskip}{2pt}\everypar{\hangindent=0.7em\hangafter=1}}X}@{}}
\toprule
\textbf{Step} & \hdr{ggmCode}{Ordinary code} & \hdr{ggmAgent}{Agent} & \hdr{ggmAnalyst}{Analyst} \\
\midrule
Set up the scenario
  & Apply edits \par Compare workbooks
  & Find the rows a storyline touches \par Propose edits and knock-on choices \par Propose a node structure
  & Set the question, storyline and node structure \par Approve edits \\
\addlinespace[5pt]
Get source data
  & Download and parse databases
  & Extract values from reports \par Record source and confidence
  & Choose the anchoring outlook \par Review by confidence \\
\addlinespace[5pt]
Prepare the data
  & Convert, aggregate, split, interpolate \par Check consistency
  & Map sources to nodes \par Draft projections and demand parameters \par Explain warnings
  & Review mappings \par Decide projections and demand parameters \\
\addlinespace[5pt]
Calibrate
  & Report deviations \par Apply changes within bounds \par Run and log
  & Diagnose deviations \par Propose changes with reasons \par Research starting prices
  & Accept the calibration \par Choose between fits \par Set market power \\
\addlinespace[5pt]
Run the model
  & Run the scenario matrix \par Record inputs and solver status
  & Explain a failed run
  & \textit{Nothing routinely} \\
\addlinespace[5pt]
Read the results
  & Compare scenarios \par Compute balances
  & Draft explanations \par Flag implausible patterns
  & Interpret \par Decide what to claim \\
\bottomrule
\end{tabularx}
\end{table}

\begin{figure}[htbp]
\centering
% Figure 5.1: the main parts of the workflow placed by judgment needed and by how quietly an error passes.
\begin{tikzpicture}[ggm, x=1.07cm, y=1.22cm]
  % zones: one band per actor, touching but not overlapping
  \begin{scope}[on background layer]
    \fill[ggmCode!11, rounded corners=4mm] (0,0) rectangle (5.0,2.0);
    \fill[ggmAgent!10, rounded corners=4mm] (1.3,2.15) rectangle (10.25,4.35);
    \fill[ggmAnalyst!13, rounded corners=4mm] (4.9,4.5) rectangle (10.25,7.0);
  \end{scope}
  \node[font=\sffamily\footnotesize\bfseries, text=ggmCode, anchor=south west, inner sep=1.8mm] at (0,0) {ORDINARY CODE};
  \node[font=\sffamily\footnotesize\bfseries, text=ggmAgent, anchor=north west, inner sep=1.8mm] at (1.3,4.35) {AGENT PROPOSES, CODE CHECKS};
  \node[font=\sffamily\footnotesize\bfseries, text=ggmAnalyst!85!black, anchor=north east, inner sep=1.8mm] at (10.25,7.0) {ANALYST DECIDES};
  % axes
  \draw[ggmMuted, line width=0.7pt, -{Stealth[length=2.2mm]}] (-0.2,-0.2) -- (10.45,-0.2);
  \draw[ggmMuted, line width=0.7pt, -{Stealth[length=2.2mm]}] (-0.2,-0.2) -- (-0.2,7.2);
  \node[font=\sffamily\scriptsize, text=ggmMuted, anchor=north west] at (-0.2,-0.26) {error is loud: shows at once};
  \node[font=\sffamily\scriptsize, text=ggmMuted, anchor=north east] at (10.45,-0.26) {error is quiet: shows only in the conclusions};
  \node[font=\sffamily\small\bfseries, text=ggmInk, anchor=north] at (5.1,-0.62) {How quietly an error passes};
  \node[font=\sffamily\small\bfseries, text=ggmInk, rotate=90, anchor=south] at (-0.68,3.5) {Judgment needed};
  \node[font=\sffamily\scriptsize, text=ggmMuted, rotate=90, anchor=south west] at (-0.25,-0.2) {rule-bound};
  \node[font=\sffamily\scriptsize, text=ggmMuted, rotate=90, anchor=south east] at (-0.25,7.2) {a view of the market};
  % points: \pt{x}{y}{colour}{label}
  \newcommand{\pt}[4]{\filldraw[fill=#3, draw=white, line width=0.9pt] (#1,#2) circle (0.12);
    \node[font=\sffamily\footnotesize, text=ggmInk, anchor=west, inner sep=2.4pt, align=left] at (#1+0.05,#2) {#4};}
  % code
  \pt{0.35}{0.8}{ggmCode}{Convert\\units}
  \pt{0.35}{1.5}{ggmCode}{Run the\\model}
  \pt{2.35}{0.8}{ggmCode}{Compute shipping\\distances}
  \pt{2.35}{1.5}{ggmCode}{Check workbook\\consistency}
  % agent
  \pt{1.6}{2.75}{ggmAgent}{Extract values\\from reports}
  \pt{4.35}{2.75}{ggmAgent}{Map sources\\to nodes}
  \pt{7.1}{2.75}{ggmAgent}{Draft explanations\\of results}
  \pt{1.6}{3.55}{ggmAgent}{Explain load\\warnings}
  \pt{4.35}{3.55}{ggmAgent}{Find the rows a\\storyline touches}
  \pt{7.1}{3.55}{ggmAgent}{Diagnose\\deviations}
  % analyst
  \pt{5.2}{4.95}{ggmAnalyst}{Market\\power}
  \pt{7.7}{4.95}{ggmAnalyst}{Choose between\\equal fits}
  \pt{5.2}{5.7}{ggmAnalyst}{Storyline and\\projections}
  \pt{7.7}{5.7}{ggmAnalyst}{Accept the\\calibration}
  \pt{5.2}{6.35}{ggmAnalyst}{Interpret results}
\end{tikzpicture}

\caption{The main parts of the workflow, placed by how much judgment they need and by how quietly an error in them would pass unnoticed. Colours show where each part belongs.}
\label{fig:task-map}
\end{figure}

Four findings carry into the design.
First, ordinary code has a part in every step, and running the model needs nothing else.
Second, the agent's part is reading, mapping, diagnosing and drafting, never computing: wherever a rule gives one correct answer, code should give it.
Third, the judgment lies in two places, as Section 3.5 found: the storyline with its projections, and the calibration.
Neither has a correct answer to check against, so an agent can prepare the material but not decide.
Fourth, outside these two places the person's role changes rather than disappears.
In the data steps, the analyst moves from doing the work to reviewing it, which is only an improvement if the review is quicker than the work and catches the errors that matter.
Much of the design that follows is about making that review quick and reliable.
